\section{Introduction}\label{sec:introduction}

Bose--Einstein condensation is a statement about coherence across a
macroscopic sample. For an interacting gas, local density estimates or
accurate thermodynamic energies do not by themselves give that coherence.
The question addressed here is whether a dilute three-dimensional gas
with an impenetrable core has a macroscopically occupied orbital at a
positive temperature that stays fixed as the volume grows.

\subsection{The model and the result}

Fix an exclusion distance $a>0$. On the flat torus
$\Lambda_L=(\R/L\Z)^3$, let $d_L$ denote Euclidean torus distance and set
\[
 \Omega_{N,L,a}
 =\{(x_1,\ldots,x_N)\in\Lambda_L^N:
                   d_L(x_i,x_j)>a\text{ whenever }i\ne j\}.
\]
In units $\hbar^2/(2m)=k_B=1$, the hard-sphere Hamiltonian
$H_{N,L,a}$ is the self-adjoint operator associated with
\[
 q[f]=\sum_{i=1}^N\int_{\Omega_{N,L,a}}|\nabla_i f|^2,
 \qquad f\in H^1_0(\Omega_{N,L,a})\text{ symmetric}.
\]
For $T>0$, its canonical state and one-particle density matrix are
\[
 \Gamma_{N,L,T}=
 \frac{e^{-H_{N,L,a}/T}}{\Tr e^{-H_{N,L,a}/T}},\qquad
 \gamma^{(1)}_{N,L,T}=N\Tr_{2,\ldots,N}\Gamma_{N,L,T}.
\]
We use zero extension from the allowed configuration set when taking
partial traces. Thus $\gamma^{(1)}_{N,L,T}$ acts on $L^2(\Lambda_L)$
and has trace $N$. The constant orbital is $u_{0,L}=L^{-3/2}$.
Macroscopic occupation of this orbital means that
$N^{-1}\langle u_{0,L},\gamma^{(1)}_{N,L,T}u_{0,L}\rangle$ has a
positive lower limit as $L\to\infty$ with $N/L^3$ tending to a fixed
density. Translation invariance makes the constant orbital an eigenvector
of the density matrix, so this gives condensation in the sense of
\citet{PenroseOnsager1956}.

\begin{theorem}\label{thm:main}
For every $a>0$, there is $\rho_*(a)>0$ such that for each fixed
$0<\rho<\rho_*(a)$ there is a fixed temperature $T=T(a,\rho)>0$ for which
\[
 \liminf_{\substack{L\to\infty\\N/L^3\to\rho}}
 \frac{\langle u_{0,L},\gamma^{(1)}_{N,L,T}u_{0,L}\rangle}{N}>0.
\]
The exclusion distance, density, and temperature are held fixed in this
limit.
\end{theorem}

The conclusion concerns the exact canonical Gibbs state. The temperature
obtained by the proof is small and is not intended to locate the phase
transition. Its role is to provide a strictly positive temperature at
every sufficiently small fixed gas parameter $\rho a^3$.

\subsection{Historical context and related methods}

The statistical origin of Bose--Einstein condensation is Bose's counting of
light quanta and Einstein's extension to material particles
\citep{Bose1924,Einstein1924IdealGas,Einstein1925IdealGas}.
For interacting particles, Penrose and Onsager formulated condensation through
a macroscopic eigenvalue of the one-particle density matrix
\citep{PenroseOnsager1956}.  In a translation-invariant gas, the constant
orbital is the natural candidate.  Bogoliubov's theory expands about its
macroscopic occupation to obtain the low-energy excitation spectrum
\citep{Bogoliubov1947}; the hard-sphere calculation of Lee, Huang and Yang
predicts the leading energy correction and depletion
\citep{LeeHuangYang1957}.  Establishing the occupation for the exact equilibrium
state is a separate mathematical question from deriving these expansions.

For hard spheres, rigorous energy bounds began with Dyson
\citep{Dyson1957}, and Lieb and Yngvason obtained the leading dilute
asymptotics \citep{LiebYngvason1998}.  Fournais and Solovej proved the
Lee--Huang--Yang lower bound for interactions including hard cores
\citep{FournaisSolovej2023}; Basti, Brooks, Cenatiempo, Olgiati and Schlein
proved the matching upper bound for hard spheres and then for general
nonnegative short-range interactions
\citep{BastiEtAl2026HardSpheres,BastiEtAl2026GeneralPotentials}.
At positive temperature, Fournais, Junge, Girardot, Morin, Olivieri and Triay
obtained a low-temperature free-energy lower bound that includes hard cores
\citep{FournaisEtAl2026FreeEnergy}.  These results control thermodynamic
energies.  They do not by themselves give a volume-uniform lower bound on a
single orbital's occupation: the one-particle kinetic gap tends to zero as
the box grows.

Rigorous continuum condensation has been established in several scaling
regimes.  Lieb and Seiringer proved complete ground-state condensation in the
trapped Gross--Pitaevskii limit \citep{LiebSeiringer2002}.  At positive
temperature, Deuchert, Seiringer and Yngvason treated a trapped gas in a
combined thermodynamic and Gross--Pitaevskii limit
\citep{DeuchertSeiringerYngvason2019}, and Deuchert and Seiringer treated the
homogeneous canonical gas in Gross--Pitaevskii scaling
\citep{DeuchertSeiringer2020GP}.  The latter theorem allows hard cores, but
their scattering length is proportional to $L/N$ and therefore vanishes at
fixed density as $N$ grows.  Further results of Fournais, Chong, Liang and Nam,
and Junge allow increasing box sizes coupled to increasing dilution
\citep{Fournais2021,ChongLiangNam2026,Junge2026}.  Such coupled limits do not
permit unrestricted volume growth at a fixed positive gas parameter.
Deuchert and Seiringer also established positive-temperature condensation in
a combined semiclassical and mean-field limit for a trapped gas
\citep{DeuchertSeiringer2021MeanField}.

Fixed-parameter thermodynamic condensation also has precedents with different
interactions and state formulations.  Aizenman, Lieb, Seiringer, Solovej and
Yngvason proved low-temperature condensation for a hard-core lattice gas at
half filling \citep{AizenmanLiebSeiringerSolovejYngvason2004}.  S\"ut\H{o}'s
continuum theorem assumes both the potential and its Fourier transform to be
nonnegative, with integrability and regularity hypotheses that exclude a
literal hard core \citep[Theorem~1.2]{Suto2023TwoStep}.  Galanda and Pinamonti
construct algebraic equilibrium states for smooth interactions in coupled
weak-interaction limits and, after renormalization, in Gross--Pitaevskii
scaling with vanishing temperature
\citep{GalandaPinamonti2025Equilibrium,GalandaPinamonti2026GP}.
The present theorem concerns the exact canonical hard-sphere state with the
exclusion distance, positive density, and positive temperature all fixed
before the volume grows.  It gives a positive condensate fraction at
sufficiently low temperature; it does not determine the transition
temperature or the leading depletion.

Two companion articles provide the immediate starting point.
\emph{Ground-state condensation in the dilute hard-sphere gas}
\citep{HS} treats the fixed-parameter ground-state problem and develops
hard-core bridge replacements, movable groups of particles, routes avoiding
unusable cells, and cancellation of conditional normalizers.
\emph{A density-uniform condensate bound for dilute Bose gases}
\citep{B} treats bounded finite-range repulsive potentials at positive
temperature and supplies localization and short-slab trajectory estimates,
as well as the underlying lattice-path construction.  We reproduce the
arguments needed here.  The additional issue is to obtain joint spatial and
trajectory estimates for the hard-core canonical Gibbs state, and to make
the short-path comparison control its one-particle density matrix.

The local analysis builds on established tools.  The hard-core ray estimate
belongs to the Dyson local-energy method, and the subdivision and
two-particle bounds are related to the work of Lieb and Yngvason and of
Lundholm, Portmann and Solovej
\citep{Dyson1957,LiebYngvason1998,LundholmPortmannSolovej2015}.
Particle insertion has antecedents in the nearest-neighbor trial states of
Dyson and Lieb, Seiringer and Yngvason
\citep{Dyson1957,LiebSeiringerYngvason2000}.
Localization into two particle species uses the Fock-space localization
formalism of Derezi\'nski and G\'erard, its many-body development by Lewin,
and its thermal application by Deuchert, Seiringer and Yngvason
\citep{DerezinskiGerard1999,Lewin2011,DeuchertSeiringerYngvason2019}.
Strong subadditivity and monotonicity of relative entropy provide the
entropy comparisons \citep{LiebRuskai1973Entropy,Lindblad1975}.
The use of positive imaginary-time paths goes back to Feynman's formulation
\citep{Feynman1953AtomicLambda}; Ueltschi's account distinguishes cycle
information from one-particle off-diagonal order
\citep{Ueltschi2006FeynmanCycles}.  Here the overlap identity addresses the
one-particle occupation directly.

The long-distance comparison has a separate probabilistic ancestry.
Benjamini, Pemantle and Peres constructed paths with exponential intersection
tails \citep{BPP1998}; H\"aggstr\"om and Mossel developed a construction using
random biases on blocks at several scales \citep{HaggstromMossel1998}.
The finite-endpoint averaging method of Abbe, Massouli\'e, Montanari, Sly and
Srivastava, and its adaptation by Garban and Spencer, turn intersection
control into long-range information or order
\citep{AbbeMassoulieMontanariSlySrivastava2018,GarbanSpencer2023}.
Averaging path likelihoods before taking their second moments also appears
in trail detection and in the point-reassignment argument of Peres and Sly
\citep{AriasCastroCandesHelgasonZeitouni2008,BergerPeres2013,PeresSly2014}.
Our lattice construction follows the explicit skeleton in \citet[Section~8.1]{B}
and the obstacle detours in \citet[Section~7]{HS}, using Tim\'ar's
boundary-connectivity theorem \citep[Theorem~4]{Timar2013}.
The hard-core comparison retains the exact interaction normalizers; after
averaging two route changes, their cost is controlled by mutual encounters
rather than by the length of either route.


\subsection{The argument}

The starting point is the short-path comparison developed for the
hard-sphere ground state in \citet{HS}. Its use at positive temperature
requires two changes. First, the local hypotheses must hold jointly for
arbitrarily many prescribed cells under a thermal path measure. Second,
the path comparison must control the one-particle density matrix of a
mixed state. We prove the needed thermal estimates and retain the
deterministic comparison mechanism. The bridge, lattice, and measure
comparison arguments used below are included with proofs.

To identify the overlap we must bound, let $\Phi(X,W')$ be the
nonnegative integral kernel of $\Gamma_{N,L,T}^{1/2}$, with
$X,W'\in\Lambda_L^N$. Write $X=(x,X')$, where $X'$ is the ordered
configuration of the other $N-1$ particles, called the bath. Factoring
the Gibbs kernel through its square root gives
\[
 \frac{\langle u_{0,L},\gamma^{(1)}_{N,L,T}u_{0,L}\rangle}{N}
 =\frac1{L^3}\int
 \Phi(x,X',W')\Phi(y,X',W')\,
 \dd x\,\dd y\,\dd X'\,\dd W'.
\]
Thus the two factors have the same bath $X'$ and retained configuration
$W'$, while the distinguished particle moves from $x$ to $y$.
Section~\ref{sec:slab} realizes $\Phi(X,W')^2$ as the joint law of
the midpoint configuration and $W'$ on a short Brownian time interval.
Once the interval endpoints and $W'$ are fixed, the conditional law is a
product of Brownian bridge laws conditioned on hard avoidance. This representation lets
us analyze local changes against deterministic obstacle paths while
retaining the mixed-state overlap above.

The estimates are made in the scaled units of Section~\ref{sec:setup},
in which a unit cell contains $D\gg1$ particles and the core radius is
$\alpha/D$, with $\alpha$ in a fixed interval.
Sections~\ref{sec:energy} and~\ref{sec:spatial} give joint bounds for
low cell counts and excessively close neighbors. These estimates use
entropy as well as energy: a relative-entropy inequality controls a
diagonal change of density, and a family of normalized insertion maps
with a uniform frame bound controls the free-energy cost of adding a
particle, even in bounded attractive fields.

The remaining local conditions concern whole trajectories.
Section~\ref{sec:slab} marks paths that violate a condition, assigns
auxiliary particle species to the marks, and removes their hard-core
constraints in an upper bound. A trace estimate for a heat semigroup
between two Lipschitz multipliers converts the marked path moments into
free-energy estimates; exterior powers permit control of every positive
trace power. These arguments adapt the thermal methods of \citet{B} to
Dirichlet hard-core domains. They yield an exponential bound in the number
of prescribed defective cells. This joint bound will control obstructing
clusters when routes are drawn through the remaining cells.

The local bridge estimates of Section~\ref{sec:bridges} show that most
cells contain many labels whose paths can be moved through any adjacent
cell with uniformly positive success probability.
Section~\ref{sec:groups} splits these labels into groups and exposes all
paths outside one group. The usable cells are then determined by this
exposed environment. Section~\ref{sec:lattice} constructs random routes
between many pairs of cells, with a bounded exponential moment for the
encounters of two independently drawn routes in the same environment.

Section~\ref{sec:comparison} uses one such route to construct two path
configurations with shifted midpoint assignments. Removing a different
tagged particle from each leaves the same bath midpoints and the same
$W'$, as required by the overlap identity. To compare the resulting
measure with the two original midpoint laws, we average over routes before
taking the second moment of its density. Conditional normalizers cancel
away from encounters between two routes, so the cost is controlled by
their encounters rather than by their lengths. The lattice estimate bounds
that cost, and a common-measure inequality gives positive overlap for a
positive fraction of all cell pairs.

The order of the choices is essential. The path duration and the group
size are fixed before the density threshold, and all are fixed before
the volume tends to infinity. Section~\ref{sec:physical} returns to
physical units and verifies directly that the resulting temperature is
independent of that volume.
